\documentclass[11pt]{amsart}
\usepackage[margin=1.15in]{geometry}
\usepackage[T1]{fontenc}
\usepackage{lmodern,microtype}
\usepackage{amsmath,amssymb,amsthm,mathtools}
\usepackage[colorlinks=true,linkcolor=blue,citecolor=blue,urlcolor=blue]{hyperref}

\newtheorem{theorem}{Theorem}[section]
\newtheorem{lemma}[theorem]{Lemma}
\newtheorem{corollary}[theorem]{Corollary}
\newtheorem{conjecture}[theorem]{Conjecture}
\theoremstyle{remark}
\newtheorem{remark}[theorem]{Remark}
\numberwithin{equation}{section}

\newcommand{\up}{\operatorname{up}}

\hypersetup{pdftitle={Single-lamination cluster polynomials of a disk are real-rooted},pdfauthor={Brandon Li}}

\title{Single-lamination cluster polynomials of a disk are real-rooted}
\author{Brandon Li}
\address{Great Neck, New York, USA}
\email{brandon.li.gn@gmail.com}
\subjclass[2020]{13F60, 05A20, 05E99}
\keywords{cluster algebra, lamination, shear coordinates, log-concavity, real-rooted polynomial, fence poset}

\begin{document}

\begin{abstract}
For a cluster algebra of a triangulated polygon with coefficients given by a
single lamination, Kang, Lee and Lim specialize the cluster variable of an arc
by setting every initial cluster variable to $1$ and the coefficient variable to
$q$. They proved that the resulting polynomial is unimodal and conjectured that
it is log-concave. We prove that it is either constant or has only simple,
negative real roots. In particular it is log-concave, which proves their
conjecture. The proof reduces the polynomial to a weighted count of walks with
two states, which we write as $\det(I+qMM^{\mathsf T})$ for a matrix $M$ whose
inverse is bidiagonal.
\end{abstract}

\maketitle

\section{Introduction}

Let $(S,M)$ be an unpunctured disk with $n\ge4$ marked points on its boundary,
that is, a convex $n$-gon, and let $T$ be a triangulation of it. Fomin,
Shapiro and Thurston~\cite{FST} associate a cluster algebra with $T$ in which
the cluster variables correspond to the arcs (diagonals) of the polygon.
Coefficients can be specified geometrically by a \emph{lamination}, a
collection of non-crossing curves on the surface: by Fomin and
Thurston~\cite{FT}, the shear coordinates of a lamination with respect to $T$
can be used as frozen rows of the exchange matrix, and the resulting extended
exchange matrix mutates compatibly with flips of $T$.

Kang, Lee and Lim~\cite{KLL} studied the case of a \emph{single lamination}
$L$, consisting of one curve whose two endpoints lie on the boundary of the
disk. The extended exchange matrix then has one frozen row, so the cluster
variable $x_\gamma$ of an arc $\gamma$ is a Laurent polynomial in the initial
cluster variables and one coefficient variable $y$. Following~\cite[Definition~20]{KLL},
the \emph{$c$-polynomial} of $\gamma$ is
\[
c_\gamma(q)=x_\gamma\big|_{x_1=\cdots=x_{n-3}=1,\ y=q}.
\]
It has nonnegative integer coefficients. Kang, Lee and Lim proved that
$c_\gamma$ is unimodal~\cite[Theorem~14]{KLL} and stated the following
conjecture.

\begin{conjecture}[{\cite[Conjecture~1]{KLL}}]\label{conj:KLL}
Let $(S,M)$ be an unpunctured disk and let $L$ be a single lamination. For any
arc $\gamma$ on $(S,M)$, the coefficient sequence of $c_\gamma(q)$ is
log-concave.
\end{conjecture}

They point out that the standard method of proving log-concavity by geometric
lattices does not apply here: the polynomial can have constant and leading
coefficient both equal to $2$. Our main result is the following.

\begin{theorem}\label{thm:main}
Let $T$ be a triangulation of an unpunctured disk, $L$ a single lamination and
$\gamma$ an arc. Then $c_\gamma(q)$ is either a positive constant or has only
simple, negative real roots.
\end{theorem}

A polynomial with positive coefficients and only real roots is log-concave
with no internal zeros by Newton's inequalities, and hence unimodal. Thus
Theorem~\ref{thm:main} proves Conjecture~\ref{conj:KLL} and gives a second
proof of~\cite[Theorem~14]{KLL}. Since the specialization is a ring
homomorphism, we also obtain the following.

\begin{corollary}\label{cor:monomials}
With $T$ and $L$ as in Theorem~\ref{thm:main}, the specialization of every
cluster monomial at $x_1=\cdots=x_{n-3}=1$, $y=q$ has only negative real roots.
\end{corollary}

The roots of a cluster monomial may be repeated, for instance for the square of
a cluster variable.

The proof has two steps. In Section~\ref{sec:walks} we show that a weighted
count of walks on two states, in which each step allows only one kind of
change, is a determinant $\det(I+qMM^{\mathsf T})$, where $M^{-1}$ is
bidiagonal. This forces the roots to be real, negative and simple. In
Section~\ref{sec:lamination} we express $c_\gamma$ as such a count. The arc and
the lamination cross a common run of diagonals, and on this run the shear
coordinates telescope (Lemma~\ref{lem:shear}). The rest of the arc contributes
only two pairs of positive multiplicities, which become the weights of two
extra steps. No restriction on the triangulation is needed. Section~\ref{sec:examples}
compares the theorem with the examples of~\cite{KLL} and describes the
computer checks in the accompanying repository.

The hypothesis that the surface is a disk is used through the fact that the
dual graph of a triangulation of a disk is a tree. Til\`epe~\cite{Til} gives
exchange matrices of annular type with one frozen row whose specializations
are not log-concave, such as $1+q^2$, and discusses their realization by
single laminations of an annulus.

\section{Walks with two states}\label{sec:walks}

A \emph{step list} is a finite sequence of steps $1,\dots,m$, where step $j$
has a \emph{direction} $\delta_j\in\{0,1\}$ and a \emph{weight} $w_j>0$. A
\emph{walk} is a sequence $(z_0,z_1,\dots,z_m)\in\{0,1\}^{m+1}$ with
$z_j\in\{z_{j-1},\delta_j\}$ for every $j$: at step $j$ the walk either stays
or moves to the state $\delta_j$. Step $j$ is a \emph{change} if
$z_j\ne z_{j-1}$ and an \emph{upward change} if moreover $z_j=1$. The weight of
a walk is the product of $w_j$ over its changes, and $\up(z)$ denotes its
number of upward changes. For $\sigma,\tau\in\{0,1\}$ let
\[
W_{\sigma\tau}(q)=\sum_{z}\Bigl(\prod_{j\text{ a change}}w_j\Bigr)q^{\up(z)},
\]
the sum over walks with $z_0=\sigma$ and $z_m=\tau$. For a nonzero Laurent
polynomial $P$ in $q$, write $\widetilde P$ for $P$ divided by the lowest power
of $q$ occurring in it.

\begin{lemma}\label{lem:walks}
If $W_{\sigma\tau}\ne0$, then $\widetilde W_{\sigma\tau}$ is a positive
constant or has only simple, negative real roots.
\end{lemma}

\begin{proof}
\emph{Reduction.} If steps $j$ and $j+1$ have the same direction, a walk
changes at most once at these two steps, and the state after them does not
depend on which of the two it uses. Replacing them by one step of the same
direction and weight $w_j+w_{j+1}$ therefore leaves $W_{\sigma\tau}$
unchanged. After all such merges, the directions alternate. A first step with
direction $\sigma$ can never be a change, and a last step with direction
different from $\tau$ cannot be a change in a walk ending at $\tau$; deleting
such steps also leaves $W_{\sigma\tau}$ unchanged. We may therefore assume that
the directions alternate, that $\delta_1=1-\sigma$, and that $\delta_m=\tau$
(if $m=0$, then $\sigma=\tau$ and $W_{\sigma\tau}=1$).

\emph{Equal endpoints.} Let $\sigma=\tau$, so $m=2s$. Write the weights as
$a_1,b_1,\dots,a_s,b_s$, so that the steps with weight $a_i$ move away from
$\sigma$ and those with weight $b_i$ move back. A walk with $r$ excursions away
from $\sigma$ is determined by the steps $a_{i_1},b_{j_1},\dots,a_{i_r},b_{j_r}$
where it changes, and these are exactly the index sequences with
\begin{equation}\label{eq:interlace}
i_1\le j_1<i_2\le j_2<\cdots<i_r\le j_r.
\end{equation}
Each excursion contains exactly one upward change. Hence the coefficient of
$q^r$ in $W_{\sigma\sigma}$ is the sum of $a_{i_1}\cdots a_{i_r}b_{j_1}\cdots b_{j_r}$
over all sequences~\eqref{eq:interlace}.

Let $U$ be the $s\times s$ matrix with $U_{ij}=1$ for $i\le j$ and $U_{ij}=0$
otherwise, and let $A=\operatorname{diag}(a_i)$ and $B=\operatorname{diag}(b_i)$.
We claim that
\begin{equation}\label{eq:det}
W_{\sigma\sigma}(q)=\det(I+qAUBU^{\mathsf T}).
\end{equation}
The coefficient of $q^r$ on the right is the sum of the principal $r\times r$
minors of $AUBU^{\mathsf T}$. By the Cauchy--Binet formula, the principal minor
on a row set $I$ equals $\sum_J a_I\,b_J\,\det(U_{I,J})^2$, where $J$ runs over
the $r$-subsets of $\{1,\dots,s\}$, $a_I=\prod_{i\in I}a_i$ and
$b_J=\prod_{j\in J}b_j$. Let $I=\{i_1<\cdots<i_r\}$ and $J=\{j_1<\cdots<j_r\}$.
Row $\ell$ of $U_{I,J}$ has a one in column $k$ exactly when $j_k\ge i_\ell$,
so the rows are indicator vectors of nested final segments of the columns. Such
a matrix has nonzero determinant only if the $\ell$-th row begins exactly at
column $\ell$ for every $\ell$, that is, $j_{\ell-1}<i_\ell\le j_\ell$; it is
then upper unitriangular. Thus $\det(U_{I,J})^2$ is $1$ for the sequences
in~\eqref{eq:interlace} and $0$ otherwise, which proves~\eqref{eq:det}.

Put $M=A^{1/2}UB^{1/2}$. Then $AUBU^{\mathsf T}=A^{1/2}(MM^{\mathsf T})A^{-1/2}$,
so $W_{\sigma\sigma}(q)=\det(I+qMM^{\mathsf T})=\prod_k(1+q\lambda_k)$, where
$\lambda_1,\dots,\lambda_s$ are the eigenvalues of the positive definite matrix
$MM^{\mathsf T}$. The roots are the numbers $-1/\lambda_k<0$. They are simple
because $K=(MM^{\mathsf T})^{-1}=N^{\mathsf T}N$, where
$N=M^{-1}=B^{-1/2}U^{-1}A^{-1/2}$ is upper bidiagonal:
\[
N_{ii}=(a_ib_i)^{-1/2},\qquad N_{i,i+1}=-(b_ia_{i+1})^{-1/2}.
\]
Hence $K$ is a real symmetric tridiagonal matrix with
$K_{i,i+1}=N_{ii}N_{i,i+1}\ne0$, and such a matrix has simple eigenvalues: in
the equation $Kv=\mu v$, the first coordinate of $v$ determines all the
others. The eigenvalues of $K$ are the numbers $1/\lambda_k$.

\emph{Different endpoints.} Let $\sigma\ne\tau$, so $m=2s-1$ with weights
$a_1,b_1,\dots,a_{s-1},b_{s-1},a_s$; again steps $a_i$ move away from $\sigma$.
Append a last step of direction $\sigma$ and weight $t>0$, called $b_s$, and
let $P_t(q)$ be the resulting enumerator of walks from $\sigma$ to $\sigma$.
Each term of $P_t$ contains $t$ at most once, and the terms containing $t$
correspond to the walks counted by $W_{\sigma\tau}$, followed by a final change
back to $\sigma$. This final change is upward exactly when $\sigma=1$, while
every excursion has one upward change. So $P_t=P_0+t\,q^{\sigma}W_{\sigma\tau}$.

Let $B_t=\operatorname{diag}(b_1,\dots,b_{s-1},t)$ and
$M_t=A^{1/2}UB_t^{1/2}$. By the first case,
$P_t(q)=\det(I+qM_tM_t^{\mathsf T})$. With $K_t=(M_tM_t^{\mathsf T})^{-1}$
and $c=a_1\cdots a_s\,b_1\cdots b_{s-1}$,
\[
\frac{P_t(q)}{t}=\frac{\det(M_tM_t^{\mathsf T})}{t}\det(K_t+qI)=c\,\det(K_t+qI).
\]
The matrix $N_t=M_t^{-1}$ is as above, and only its last row depends on $t$.
That row has the single nonzero entry $(a_st)^{-1/2}$. Since
$K_t=\sum_k n_k^{\mathsf T}n_k$ over the rows $n_k$ of $N_t$, letting
$t\to\infty$ gives
\[
q^{\sigma}W_{\sigma\tau}(q)=\lim_{t\to\infty}\frac{P_t(q)}{t}=c\,\det(K_\infty+qI),
\qquad K_\infty=\sum_{k=1}^{s-1}n_k^{\mathsf T}n_k.
\]
The rows $n_1,\dots,n_{s-1}$ are linearly independent, so $K_\infty$ is
positive semidefinite of rank $s-1$. It is tridiagonal with
$(K_\infty)_{i,i+1}=N_{ii}N_{i,i+1}\ne0$ for $1\le i\le s-1$, so its
eigenvalues are simple: $0$ and $s-1$ positive numbers $\mu_1,\dots,\mu_{s-1}$.
Hence $q^{\sigma}W_{\sigma\tau}(q)=c\,q\prod_k(q+\mu_k)$, and
$\widetilde W_{\sigma\tau}=c\prod_k(q+\mu_k)$. For $s=1$ this is the constant
$c=a_1$.
\end{proof}

\section{From a lamination to walks}\label{sec:lamination}

\subsection{Setting}
We draw the disk as a convex $n$-gon and $T$ as a set of $n-3$ diagonals
$\tau_1,\dots,\tau_{n-3}$. The single lamination $L$ is a curve whose endpoints
lie in the interiors of two boundary edges. It is not allowed to cut off a
region containing at most one marked point, so its endpoints lie on two
non-adjacent edges, and up to isotopy it is the chord joining their
midpoints. We use the conventions of~\cite[Section~2]{KLL}: the exchange
matrix $B(T)=(b_{ij})$ has $b_{ij}=1$ when $\tau_j$ immediately follows
$\tau_i$ clockwise in a triangle of $T$ (summed over the triangles), and the
shear coordinate $b_{\tau}(T,L)$ is $+1$ or $-1$ when $L$ crosses the
quadrilateral formed by the two triangles adjacent to $\tau$ in the S-shape or
the Z-shape of~\cite[Figure~3]{KLL}, and $0$ otherwise. The extended exchange
matrix is $B(T)$ with the extra row $(b_{\tau_j}(T,L))_j$.

If $\gamma$ is a diagonal of $T$ then $c_\gamma=1$, so let $\gamma$ be any
other diagonal. Orient $\gamma$, and let $\rho_1,\dots,\rho_d$ be the diagonals
of $T$ that it crosses, in order. For $1\le j<d$ the diagonals $\rho_j$ and
$\rho_{j+1}$ are two sides of a triangle of $T$; let $v_j$ be their common
endpoint. Every vertex other than the endpoints of $\gamma$ lies to the left
or to the right of $\gamma$.

\subsection{The \texorpdfstring{$c$}{c}-polynomial as an ideal sum}
The \emph{fence poset} of $\gamma$ is the partial order on $\{1,\dots,d\}$
generated by the relations $j+1\prec j$ when $v_j$ lies to the right of
$\gamma$, and $j\prec j+1$ when it lies to the left. By the expansion formulas
of Schiffler~\cite{Sch} and of Musiker, Schiffler and Williams~\cite{MSW}, in
the poset form of~\cite[Section~3.2]{KLL} (see also~\cite{PRS}), the
$F$-polynomial of $x_\gamma$ is
\[
F_\gamma(y_1,\dots,y_{n-3})=\sum_{I}\prod_{j\in I}y_{\rho_j},
\]
where $I$ runs over the order ideals of the fence poset and $y_\tau$ denotes
the principal coefficient attached to $\tau$. We encode an order ideal by
$z_j=\mathbf 1_{j\in I}$. The binary sequences $(z_1,\dots,z_d)$ that arise are
exactly those in which, for each $j<d$, a change from $z_j=0$ to $z_{j+1}=1$
is allowed only if $v_j$ lies to the right of $\gamma$, and a change from $1$
to $0$ only if $v_j$ lies to the left.

Now specialize. The coefficients $y_\tau=q^{b_\tau(T,L)}$ lie in the tropical
semifield generated by $q$. By the separation formula of Fomin and
Zelevinsky~\cite[Theorem~3.7]{FZ},
$x_\gamma=F_\gamma(\hat y)\,/\,F_\gamma|_{\mathrm{Trop}}(y)$ times a Laurent
monomial in the initial cluster variables. When all initial cluster variables
equal $1$, we have $\hat y_\tau=y_\tau$, and the tropical evaluation of the
subtraction-free polynomial $F_\gamma$ is $q$ raised to the smallest exponent
occurring in $F_\gamma(q^{b_\tau(T,L)})$. Therefore
\begin{equation}\label{eq:idealsum}
c_\gamma(q)=\widetilde{E}(q),\qquad
E(q)=\sum_{z}q^{\sum_{j=1}^d\beta_jz_j},\qquad \beta_j=b_{\rho_j}(T,L),
\end{equation}
where $z$ runs over the binary sequences described above. Reversing the
orientation of the fence relations, or the sign of every shear coordinate,
replaces $c_\gamma(q)$ by $q^{D}c_\gamma(1/q)$ with $D=\deg c_\gamma$. This
does not affect the conclusion of Theorem~\ref{thm:main}. The computations in
Section~\ref{sec:examples} confirm~\eqref{eq:idealsum} with the conventions
above against direct mutation in all polygons with at most nine vertices.

\subsection{The common run and the shear coordinates}
The triangles of $T$ are the vertices of its dual graph, which is a tree since
the disk is simply connected. The triangles met by $\gamma$ form a path in
this tree, and so do the triangles met by $L$. Two paths in a tree meet in a
path, possibly empty, so the diagonals crossed by both $\gamma$ and $L$ are
consecutive in the order along $\gamma$, say $\rho_{p+1},\dots,\rho_{p+r}$.
Only these diagonals can have nonzero shear coordinate among
$\rho_1,\dots,\rho_d$, since a diagonal that $L$ does not cross has shear
coordinate~$0$. If $r=0$, then all exponents in~\eqref{eq:idealsum} vanish and
$c_\gamma$ is a positive constant. Assume $r\ge1$, and orient $L$ so that it
crosses $\rho_{p+1},\dots,\rho_{p+r}$ in this order.

For $1\le i<r$, the diagonals $\rho_{p+i}$ and $\rho_{p+i+1}$ are sides of a
triangle that both curves cross, entering through the first and leaving
through the second. Hence their common vertex $v_{p+i}$ lies on the same side
of $L$ as of $\gamma$. Put $d_i=1$ if this vertex lies to the right of $L$, and
$d_i=0$ otherwise. Before crossing $\rho_{p+1}$, the curve $L$ passes through
the triangle on the other side of $\rho_{p+1}$, entering it through a side
$\alpha$, which is a boundary edge or another diagonal. Put $d_0=1$ if the
common endpoint of $\alpha$ and $\rho_{p+1}$ lies to the right of $L$, and
$d_0=0$ otherwise. Define $d_r$ in the same way using the side $\omega$
through which $L$ leaves the triangle after $\rho_{p+r}$. The bits $d_0$ and
$d_r$ concern $L$ only; they need not agree with the turns of $\gamma$.

\begin{lemma}\label{lem:shear}
For $1\le i\le r$, the shear coordinate of $\rho_{p+i}$ is $d_{i-1}-d_i$.
\end{lemma}

\begin{proof}
Let $Q$ be the quadrilateral formed by the two triangles adjacent to
$\rho=\rho_{p+i}$. The curve $L$ crosses $Q$, entering through a side $\alpha'$
and leaving through a side $\omega'$, where $\alpha'$ is $\rho_{p+i-1}$ (or
$\alpha$ if $i=1$) and $\omega'$ is $\rho_{p+i+1}$ (or $\omega$ if $i=r$). The
side $\alpha'$ meets $\rho$ in the endpoint of $\rho$ that bit $d_{i-1}$
describes, and $\omega'$ meets $\rho$ in the endpoint that $d_i$ describes.
Because $L$ crosses $\rho$, one endpoint of $\rho$ lies to the right of $L$
and the other to the left. If $d_{i-1}=d_i$, then $\alpha'$ and $\omega'$ share
an endpoint of $\rho$, so they are adjacent sides of $Q$ and $L$ cuts off a
corner of $Q$; the shear coordinate is $0$. If $d_{i-1}\ne d_i$, then
$\alpha'$ and $\omega'$ are opposite sides of $Q$, and $L$ forms an S-shape or
a Z-shape. In~\cite[Figure~3]{KLL}, the S-shape with contribution $+1$ is the
one in which $L$ enters through the side at the endpoint of $\rho$ on its
right and leaves through the side at the endpoint on its left, that is,
$(d_{i-1},d_i)=(1,0)$. The Z-shape is $(d_{i-1},d_i)=(0,1)$ with contribution
$-1$. Reversing the orientation of $L$ exchanges left and right and reverses
the order of the bits, so the description does not depend on the orientation.
\end{proof}

\subsection{Proof of Theorem~\ref{thm:main}}
Keep $r\ge1$. In the sum~\eqref{eq:idealsum}, the exponent depends only on
$z_{p+1},\dots,z_{p+r}$, and by Lemma~\ref{lem:shear} it equals
\[
e(z)=\sum_{i=1}^r(d_{i-1}-d_i)\,z_{p+i}.
\]
The constraints on $z_1,\dots,z_{p+1}$ and on $z_{p+r},\dots,z_d$ form two
paths, which meet the run only in $z_{p+1}$ and $z_{p+r}$. For $\kappa\in\{0,1\}$,
let $\lambda_\kappa$ be the number of admissible $(z_1,\dots,z_{p+1})$ with
$z_{p+1}=\kappa$, and $\lambda'_\kappa$ the number of admissible
$(z_{p+r},\dots,z_d)$ with $z_{p+r}=\kappa$. The constant sequences are
admissible, so $\lambda_\kappa,\lambda'_\kappa\ge1$. Inside the run, the change
between $z_{p+i}$ and $z_{p+i+1}$ is allowed only toward $d_i$, since $d_i$
records the side of $v_{p+i}$ with respect to both curves. Writing
$u_i=z_{p+i}$, we obtain
\begin{equation}\label{eq:split}
E(q)=\sum_{u}\lambda_{u_1}\,\lambda'_{u_r}\,q^{e(u)},
\end{equation}
over $u\in\{0,1\}^r$ such that each step $u_i\to u_{i+1}$ ($1\le i<r$) stays
or moves to $d_i$.

Turn~\eqref{eq:split} into a walk enumerator. Add a state $u_0=1-d_0$ and a
first step of direction $d_0$ with weight $\lambda_{d_0}/\lambda_{1-d_0}$, so
that $u_1$ is free and the step changes exactly when $u_1=d_0$. Add a state
$u_{r+1}=d_r$ and a last step of direction $d_r$ with weight
$\lambda'_{1-d_r}/\lambda'_{d_r}$, which changes exactly when $u_r=1-d_r$. Give the
steps inside the run weight $1$. The step from $u_i$ to $u_{i+1}$ has direction
$d_i$ for $0\le i\le r$, so it contributes $d_i(u_{i+1}-u_i)$ upward changes.
Using $u_0=1-d_0$ and $u_{r+1}=d_r$,
\[
\up(u)=\sum_{i=0}^{r}d_i(u_{i+1}-u_i)=d_0u_1+\sum_{i=1}^{r-1}d_i(u_{i+1}-u_i)+d_r(1-u_r)
=e(u)+d_r.
\]
Hence $E(q)=\lambda_{1-d_0}\,\lambda'_{d_r}\,q^{-d_r}\,W_{1-d_0,\,d_r}(q)$ for
this step list, and $c_\gamma=\widetilde E$ is a positive multiple of
$\widetilde W_{1-d_0,\,d_r}$. The constant sequences show that
$W_{1-d_0,d_r}\ne0$, and Lemma~\ref{lem:walks} completes the proof.
\qed

\begin{proof}[Proof of Corollary~\ref{cor:monomials}]
A cluster monomial is a product of cluster variables, and setting
$x_1=\cdots=x_{n-3}=1$, $y=q$ is a ring homomorphism. Each factor specializes
to a positive constant or a polynomial with only negative roots, by
Theorem~\ref{thm:main}.
\end{proof}

\section{Examples and computations}\label{sec:examples}

\subsection{The examples of Kang, Lee and Lim}
The worked examples in~\cite[Section~5.1]{KLL} (numbered as in
arXiv:2508.04396v4) print some $c$-polynomials that do not have only real
roots. The following comparison shows that they do not
contradict Theorem~\ref{thm:main}.

In~\cite[Example~3]{KLL} the extended exchange matrix is the one printed
in~\cite[Example~2]{KLL}, for a triangulated octagon, and the printed value is
$2q^2+2q+2$. In~\cite[Example~4]{KLL} the printed value is $2q^3+4q^2+2q+1$.
Mutating the printed matrices along the five diagonals in order gives
$2q^2+6q+3$ and $2q^2+5q+2$, respectively, and no run of consecutive diagonals
gives the printed values. In fact, neither printed polynomial is the
$c$-polynomial of any arc for any single lamination of any triangulated
octagon. In~\cite[Example~6]{KLL} the displayed polynomial
$2q^8+4q^7+6q^6+10q^5+11q^4+12q^3+9q^2+5q+2$ coincides with the one listed
in~\cite[Table~4]{KLL} for a lamination with several curves, and it is not
log-concave, since $6^2<4\cdot10$. For the single lamination drawn
in~\cite[Figure~17]{KLL}, label the vertices $A,\dots,L$ clockwise, starting
from the leftmost one. The triangulation consists of $BL$, $CL$, $DL$, $DK$,
$EK$, $EJ$, $EI$, $FI$, $GI$, the arc is $AH$, and $L$ runs from the edge $AB$
to the edge $GH$. Direct mutation gives
\[
c_{AH}(q)=8q^3+26q^2+22q+5,
\]
whose roots are simple and negative.

\subsection{Computer checks}
The repository listed below contains short Python programs, using only the
standard library and exact arithmetic, that check each step of the argument on
small cases.
\begin{itemize}
\item Lemma~\ref{lem:walks} and the identity~\eqref{eq:det} are compared on
$11{,}575$ random step lists.
\item Lemma~\ref{lem:shear}: in every triangulated $n$-gon with $n\le10$ and for
every single lamination, the shear coordinates computed from the S/Z rule equal
$d_{i-1}-d_i$ on the diagonals crossed by $L$ and vanish elsewhere.
\item The identity~\eqref{eq:idealsum} and Theorem~\ref{thm:main}: for every
triangulation of an $n$-gon with $4\le n\le9$, every single lamination and
every arc ($375{,}040$ triples in all), mutation of the extended exchange matrix
gives the same polynomial as the ideal sum. Each of the $228$ distinct
polynomials that occur is checked by Sturm's theorem to have only simple,
negative roots.
\item The computations of this section.
\end{itemize}

\section*{Verification code}

The programs described in Section~\ref{sec:examples} are available at\\
\url{https://github.com/brandonlign/disk-lamination}. The script
\path{checks/run_all.sh} runs all of them in under a minute with Python~3.

\section*{Acknowledgments}
Claude (Anthropic) helped with the proofs, the verification code, and the
manuscript.

\begin{thebibliography}{9}

\bibitem{FST}
S.~Fomin, M.~Shapiro, and D.~Thurston, \emph{Cluster algebras and triangulated
surfaces. {P}art {I}: Cluster complexes}, Acta Math. \textbf{201} (2008),
83--146.

\bibitem{FT}
S.~Fomin and D.~Thurston, \emph{Cluster algebras and triangulated surfaces.
{P}art {II}: Lambda lengths}, Mem. Amer. Math. Soc. \textbf{255} (2018),
no.~1223.

\bibitem{FZ}
S.~Fomin and A.~Zelevinsky, \emph{Cluster algebras {IV}: Coefficients},
Compos. Math. \textbf{143} (2007), 112--164.

\bibitem{KLL}
W.~Kang, K.~Lee, and E.~Lim, \emph{Unimodality and cluster algebras from
surfaces}, European J. Combin. \textbf{136} (2026), 104392;
\href{https://arxiv.org/abs/2508.04396v4}{arXiv:2508.04396v4}.

\bibitem{MSW}
G.~Musiker, R.~Schiffler, and L.~Williams, \emph{Positivity for cluster
algebras from surfaces}, Adv. Math. \textbf{227} (2011), 2241--2308.

\bibitem{PRS}
V.~Pilaud, N.~Reading, and S.~Schroll, \emph{Posets for {$F$}-polynomials in
cluster algebras from surfaces}, preprint,
\href{https://arxiv.org/abs/2311.06033}{arXiv:2311.06033}.

\bibitem{Sch}
R.~Schiffler, \emph{A cluster expansion formula ({$A_n$} case)}, Electron. J.
Combin. \textbf{15} (2008), \#R64.

\bibitem{Til}
A.~Til\`epe, \emph{Annular obstructions and skew-shape block identities},
preprint, 2026, \href{https://arxiv.org/abs/2611.14506}{arXiv:2611.14506}.

\end{thebibliography}

\end{document}
